\documentclass[../comorg-part1.tex]{subfiles}

\begin{document}
\section{Introduction}
\subsection{2 Classes of computer architecture}
\begin{figure}[h]
    \centering
    \begin{subfigure}[b]{0.45\textwidth}
        \centering
        \includegraphics[width=\textwidth]{images/vn-archi.png}
        \caption{Von Neumann architecture}
    \end{subfigure}
    \hfill
    \begin{subfigure}[b]{0.45\textwidth}
        \centering
        \includegraphics[width=\textwidth]{images/harvard-archi.png}
        \caption{Harvard Architecture}
    \end{subfigure}
    
    % \caption{A collective caption for both images.}
\end{figure}

\noindent The modern day design (and most of the course) 
is based on von Neumann's architecture.
\begin{figure}[H]
    \centering
    \includegraphics[width=0.7\textwidth]{images/vn-with-cpu.png}
    \caption*{CPU interacts with Memory}
\end{figure}

\noindent The components of a microcomputer.
\begin{figure}[H]
    \centering
    \includegraphics[width=0.7\textwidth]{images/microcomputer.png}
    \caption*{Components of a Microcomputer}
\end{figure}
\begin{enumerate}
    \item The 3 main components - Processor (CPU), Main Memory (DRAM), I/O interfaces -
    are conencted by a bus structure, where binary information can be transferred in parallel.
    \item A master clock drives the synchronous operations. 
    Operations closer to the CPU core are clocked faster, those involving external components are clocked slower.
    \item Reset circuitry - Active low signal on reset pin for a substantial duration (several clock cycles) is required to reset the CPU.
\end{enumerate}

\noindent Package on Package (PoP) for smaller devices, like iPad.
\begin{figure}[H]
    \centering
    \includegraphics[width=0.5\textwidth]{images/PoP.png}
    \caption*{Package on Package}
\end{figure}
\noindent PoP is an IC packaging technique that vertically stacks and interconnects separate packages (like the CPU and memory)
 via ball grid array connections.

\noindent This saves space, minimises track length between CPU and memory and allows separate testing of memory units.

\newpage

\noindent Example Apple A4 chip
\begin{figure}[H]
    \centering
    \includegraphics[width=0.5\textwidth]{images/a4-chip.png}
    \caption*{A4 chip}
\end{figure}
\noindent PoP design where the memory is stacked on top of the processor.

\end{document}